\chapter*{The semester as one argument}
\addcontentsline{toc}{chapter}{The semester as one argument}
The schedule assumes eighteen teaching weeks with three contact hours per week, plus roughly four to six hours of individual reading and exercises. These figures are a guide for planning; individual students will need more or less time. Chapters~1--6 deliberately move slowly, because everything later depends on them. Chapters~8 and~9 form one two-week unit, as do Chapters~11 and~12: in each pair, the first chapter prepares the ideas that the long proof in the second chapter uses, so the two should be read together. Weeks~17 and~18 bring together the earlier material instead of introducing a new list of topics.

A chapter is a week's reading, meant to be divided at its section headings rather than read in one sitting. For a long proof, a single sitting might cover only the statement and one of its smaller arguments. When you come back, start by saying what has already been established and what remains to be done. The worked examples are part of the main path through each chapter, so do not skip them. In a taught course, the four core exercises form the required weekly work, and the two transfer exercises provide extra practice once the core exercises are secure.

\begingroup\small
\begin{longtable}{@{}p{0.34in}p{2.12in}p{2.87in}@{}}
\toprule Week & Mathematical development & Concrete output\\\midrule\endhead
1 & Variables, representation, and explanation & A precise claim and a fully explained parity proof.\\
2 & Logic and dependence & A theorem translated, negated, and tested.\\
3 & Sets and maps & A composition proof and a counterexample to a converse.\\
4 & Building and reading arguments & A proof annotated by goal, fact, and justification.\\
5 & Induction and invariants & A recurrence proof and a terminating algorithm.\\
6 & Equivalence and symmetry & A well-defined construction on residue classes.\\
7 & Graphs, counting, and extremal choice & A counting proof and a shortest-path argument.\\
8 & Assignments and augmenting paths & An algorithm explained with its invariant.\\
9 & Hall's theorem and proof reconstruction & A complete long proof and the midterm portfolio.\\
10 & Vectors, linearity, and representation & A kernel criterion and an inner-product inequality.\\
11 & Distance, limits, and completeness & A convergence argument with explicit tolerances.\\
12 & The contraction theorem & A full proof with a numerical error certificate.\\
13 & Averaging and finite probability & An existence proof and a deterministic construction.\\
14 & Sign sequences and research experiments & Exact correlations and a necessary-order theorem.\\
15 & Lean: statements and proof states & Small proofs reconstructed in a formal language.\\
16 & Lean: definitions, trust, and repair & A formalization with an assumption audit.\\
17 & Learning an unfamiliar proof with AI & A prerequisite map and a reading record for a research proof.\\
18 & A bounded research investigation & A complete capstone baseline and a defended extension.\\
\bottomrule
\end{longtable}
\endgroup

\section*{A repeatable weekly rhythm}
A week of class can follow a simple pattern. In the first meeting, unpack the opening problem and work out one example together. In the second, reconstruct the main proof and compare two possible next moves at a difficult step. In the third, examine a proposed argument---one written by a classmate, taken from the text, or produced by an AI assistant---and decide exactly what it proves and where it might fail. Students working on their own can use the same sequence: a first attempt, a guided reading, a reconstruction with the book closed, and a transfer exercise. If you already know an example, explain what happens to it when one assumption is changed instead of skipping it.

\section*{Assessment that rewards understanding}
One reasonable weighting is 40\% for the weekly core exercises, 25\% for the Week~9 portfolio, and 35\% for the final investigation. AI assistance is allowed as long as its role is stated openly. Every submission should contain a short explanation, written without assistance, of its central step; one boundary case that has been tested; and a statement of anything that is borrowed from elsewhere or not yet checked. The grade should reflect the mathematics and the understanding the student demonstrates, not the length of a conversation log or the polish of generated prose.

The Week~9 portfolio contains a correct proof of Hall's theorem, one example that satisfies Hall's condition, one that fails it, and an explanation of why a greedy assignment can get stuck. The final investigation contains an exact claim, a short list of the background it needs, a complete proof of a baseline result, a checked experiment, and one justified variation. Chapter~18 works through a complete model, and Appendix~\ref{app:guide} gives a grading rubric.

\section*{Keeping track of evidence}
Keep two separate columns in your notebook. The first records your understanding of a statement: \emph{can explain}, \emph{can follow}, or \emph{not yet understood}. The second records the support for it: \emph{proved here}, \emph{cited from another source}, \emph{checked by computer in a finite range}, or \emph{not established}. The two columns measure different things. A proof checked by a computer can still be poorly understood, and a pattern you understand well from experiments can still be unproved. Keeping both columns stops these two kinds of confidence from blurring into one.
